%% ===========================================================================
%%  main-zh.tex -- 中文示例报告（用 XeLaTeX 编译）
%%  构建：.\build.ps1        或   xelatex main-zh.tex 两遍
%%  这份文件同时是模板，新报告复制它，替换正文与元信息即可。
%% ===========================================================================
\documentclass[logo, zh]{boogonreport}

%% ----------------------------- 元信息 --------------------------------------
%% 只在这里写一次：标题、作者、编号、版本、链接。PDF 元数据会自动同步。
\reportsetup{
  title       = {技术报告模板：设计与用法},
  subtitle    = {boogonreport v1.0 的版面决策、排版规范与工程化},
  author      = {BoogonClothman},
  affiliation = {Boogon},
  email       = {hello@example.com},
  link        = {https://github.com/BoogonClothman/boogon-techreport-kit},
  linktext    = {github.com/BoogonClothman/boogon-techreport-kit},
  number      = {TR-2026-01},
  version     = {v1.0.0},
  date        = {2026-10-22},
  wordmark    = {BoogonClothman},
  shorttitle  = {技术报告模板：设计与用法},
  keywords    = {技术报告, LaTeX, 排版, 双语},
}

%% ----------------------------- 绘图 ----------------------------------------
\usepackage{pgfplots}
\pgfplotsset{compat=1.18}
\usetikzlibrary{positioning,arrows.meta,fit,backgrounds}

\begin{document}

\maketitle

\begin{abstract}
这份模板把一份技术报告的排版决策收进一个类文件：A4 单栏、Palatino 正文、
竖线图注、三线表与最优值标记、代码块、附录与参考文献，并且用一条命令切换
中英文。示例报告本身说明每个部件为什么这样设计，以及哪些事情这套模板
刻意不做。全部内容可用一条命令构建，产物在任意装了 TeX Live 的机器上一致。
\keywords{技术报告，LaTeX，排版规范，中英双语}
\end{abstract}

%% 首屏主结果图：一份报告的构建耗时。
%% 数据取自本机 TeX Live 2026 上 latexmk 的两次完整构建（含 bibtex），
%% 换机器后请重测并更新坐标，不要沿用旧数字。
\begin{figure}[htbp]
  \centering
  \begin{tikzpicture}
    \begin{axis}[
      ybar,
      bar width=26pt,
      width=0.86\textwidth,
      height=5.6cm,
      enlarge x limits=0.28,
      ymin=0,
      ymax=20,
      ylabel={构建耗时（秒）},
      ylabel style={font=\small},
      symbolic x coords={英文示例, 中文示例, 空白骨架},
      xtick=data,
      xticklabel style={font=\small},
      tick style={draw=none},
      ymajorgrids,
      grid style={bgrule!45},
      axis line style={bgrule},
      nodes near coords,
      nodes near coords style={font=\scriptsize, color=bgink},
    ]
      \addplot[fill=bgbrand, draw=none] coordinates
        {(英文示例,8.5) (中文示例,15.7) (空白骨架,1.3)};
    \end{axis}
  \end{tikzpicture}
  \caption{本机完整构建耗时（两遍编译 + 参考文献）。中文示例因为要走 XeLaTeX
  与 ctex，约为英文示例的两倍，但仍然在一次编译停顿之内。}
  \label{fig:hero}
\end{figure}


\newpage
\tableofcontents
\newpage

%% ----------------------------- 正文 ----------------------------------------
\section{引言}
\label{sec:intro}

技术报告的价值在于可核对：读者能顺着数字找到实验，顺着实验找到代码。排版不是
这份价值的来源，但它决定读者愿不愿意读下去，也决定作者每写一份报告要不要把
同样的决定重做一遍。

我们具体被三件事绊住过。第一，每次新开报告都要从空白文档重调页边距、字号、
图注和表格线，而调出来的东西每份都不太一样。第二，中英混排反复踩坑：中文要
XeLaTeX 才能正确断行，英文用 pdfLaTeX 最省事，而两套环境下的观感必须一致。
第三，数字写法没有统一约定，同一张表里出现 \texttt{62.4} 和 \texttt{62.40}，
读者就有理由怀疑哪一个是笔误。

这份模板针对的就是这三件事，贡献可以归纳为三点：

\begin{contributions}
  \item 一个类文件 \texttt{boogonreport.cls} 承载全部版式决策：标题块、摘要、
        题注、三线表、代码块、参考文献与附录；
  \item 一个 \texttt{zh} 选项同时切换字体管线与本地化文案，中文报告不再需要
        手写一串 \verb|\renewcommand|；
  \item 一套构建脚本与 CI，让任何一台装了 TeX Live 的机器上，一条命令就能得到
        同一份 PDF。
\end{contributions}

全文组织如下：\secref{sec:layout} 说明版面参数及其来源；\secref{sec:type}
交代中英两条字体管线；\secref{sec:figure} 给出图表与数字的写法；
\secref{sec:build} 是构建与发布流程；\secref{sec:limit} 列出这套模板不打算
解决的问题。

\section{版面}
\label{sec:layout}

版面是 A4 单栏，左右边距 25\,mm，上边距 30\,mm，下边距 32\,mm。版心宽度因此是

\begin{equation}
  W = 210\,\text{mm} - 2 \times 25\,\text{mm} = 160\,\text{mm}
  \label{eq:width}
\end{equation}

式~\eqref{eq:width} 在 11\,pt 的 Palatino 下大约对应每行 75 个字符，比
Bringhurst 推荐的 66 字符略宽\citep{bringhurst2013elements}。这是有意的：
技术报告里公式、代码和宽表格比纯散文多，版心窄一分，公式就要多折一行。

\begin{table}[htbp]
  \centering
  \small
  \caption{版面参数。数值与 \texttt{boogonreport.cls} 中的设定一一对应。}
  \label{tab:layout}
  \begin{tabular}{@{}lll@{}}
    \toprule
    项目 & 取值 & 理由 \\
    \midrule
    纸型          & A4（210 $\times$ 297\,mm） & 国内打印与 arXiv 都接受 \\
    左右边距      & 25\,mm                     & 与常见预印本一致 \\
    上 / 下边距   & 30\,mm / 32\,mm            & 给页眉页脚留位 \\
    正文字号      & 11\,pt                     & 屏幕阅读与打印都够 \\
    行距          & 单倍（Palatino 自带较大 x-height） & 不必再拉开 \\
    段首          & 不缩进，段间 0.5 行        & 块状段落利于扫读 \\
    \bottomrule
  \end{tabular}
\end{table}

段落不缩进、段间留半行空白。中文排版传统是首行缩进两字，但技术报告多在屏幕上
阅读，块状段落更利于跳读与定位。如果目标渠道要求缩进，在导言区加一句
\verb|\setlength{\parindent}{2em}| 并把 \verb|\parskip| 归零即可。

之所以选单栏，是因为双栏在屏幕上必须放大才能读，而技术报告的主要读者就在屏幕上。
代价同样明确：同样的页数，单栏装不下双栏那么多字，所以这套模板适合 8--20 页的
报告，不适合长篇学位论文。

\section{字体与双语}
\label{sec:type}

同一份观感背后是两条字体管线。英文报告走 pdfLaTeX：T1 编码加 \texttt{mathpazo}，
正文与数学都是 URW Palladio，也就是 Palatino 的克隆。中文报告走 XeLaTeX：
\texttt{ctex} 负责断行与标点，正文换成 OpenType 的 TeX Gyre Pagella，
数学仍是 Pazo\citep{ctex,texgyre}。

之所以不把两条管线合并成一条，是因为中文没有第二条路可走。pdfLaTeX 处理
CJK 断行与字体会得到一份缺字或乱断的 PDF，所以类文件在检测到
\texttt{zh} 与 pdfLaTeX 同时出现时会直接报错，并告诉你该敲哪条命令——报错
比一份看起来正常、实则掉字的 PDF 便宜得多。

两个字体名看起来分裂，其实是同一设计的两份实现：Palladio 是 Type~1，
Pagella 是它的 OpenType 重绘，肉眼几乎分不出来。区别在于字符覆盖——Pagella
带 →、≥、µ 这些字符，Palladio 不带。英文管线下
T1 编码对 ≥ ≤ ≠ ≈ ∞ 没有定义，输入这些字符会被\textbf{静默丢弃}，
所以类里给这五个字符补了数学映射；箭头和单位符号由 \texttt{textcomp} 自带，
无需重复声明。

\begin{table}[htbp]
  \centering
  \small
  \caption{两条字体管线。中文字体集另有 \texttt{windows} 与 \texttt{fandol} 两种选择。}
  \label{tab:pipeline}
  \begin{tabular}{@{}lll@{}}
    \toprule
    项目 & 英文管线 & 中文管线 \\
    \midrule
    引擎        & pdfLaTeX          & XeLaTeX \\
    文本字体    & URW Palladio（T1）& TeX Gyre Pagella \\
    数学字体    & Pazo Math         & Pazo Math \\
    中文支持    & ---               & ctex，\texttt{scheme=plain} \\
    页面编码    & T1                & TU + T1 \\
    适用场景    & 英文报告、arXiv   & 中文报告、中英混排 \\
    \bottomrule
  \end{tabular}
\end{table}

中文还有一份「汉字字体集」要选：Windows 下可以用系统宋体
（\texttt{fontset=windows}），跨平台则用 TeX Live 自带的 Fandol
（\texttt{fontset=fandol}）。后者任何机器都能编，代价是字形偏细、PDF 里
文字不易复制；前者更好看，但只在装了对应系统字体的机器上成立。CI 用前者之外的
选项保证可编译，个人机器用系统字体保证好看——这是可复现与好看之间的一次明确
取舍，而不是疏忽。

元信息也收在一处，这是模板里最小但最省事的约定：

\begin{lstlisting}[language=TeX,
  caption={\texttt{\textbackslash reportsetup} 的典型写法。所有键都可省略，省略时用默认值。},
  label={lst:setup}]
\documentclass[logo, zh]{boogonreport}
\reportsetup{
  title    = {技术报告模板：设计与用法},
  author   = {BoogonClothman},
  number   = {TR-2026-01},
  version  = {v1.0.1},
  date     = {2026-10-22},
}
\end{lstlisting}

同一处定义会同时写进页眉、页脚、PDF 元数据与页脚的报告编号，因此不存在
「封面改了、PDF 属性没改」这类不一致。

\section{图表与数字}
\label{sec:figure}

表格一律用 \texttt{booktabs} 的三线表：顶线、中线、底线，不加竖线。表题在表
上方，图题在图下方，两者都用竖线把标签和说明分开（\texttt{表 1 | ...}）。
最优值加粗（\verb|\best{}|），次优下划线（\verb|\second{}|），并在题注里
写明含义，读者不必回头翻正文。

每张图表都要在正文里被引用，否则删掉它。中文里章节用 \verb|\secref{}|，
写成「第 3 节」；图表用 \verb|\cref{}|，写成「图 1」「表 2」。这样改章节顺序
时不会留下错误的「第 3 节」。

图~\ref{fig:arch} 是仓库结构与构建链路，用 TikZ 直接画在文档里，不依赖外部
图片，改结构时改代码即可。

%% 仓库结构与构建链路示意图（TikZ，无外部图片依赖）
\begin{figure}[htbp]
  \centering
  \begin{tikzpicture}[
      font=\small,
      box/.style={draw=bgrule, rounded corners=2pt, fill=bgbrand!4,
                  align=center, inner sep=4pt, minimum height=8mm,
                  minimum width=40mm},
      core/.style={draw=bgbranddeep, rounded corners=3pt, fill=bgbrand!12,
                   align=center, inner sep=6pt, minimum height=26mm,
                   minimum width=38mm},
      outbox/.style={draw=bgrule, rounded corners=2pt, fill=black!4,
                  align=center, inner sep=4pt, minimum height=9mm,
                  minimum width=34mm},
      arr/.style={-{Latex[length=2mm]}, bgrule!90, thick},
      note/.style={font=\scriptsize, text=bgink, align=center},
      node distance=4mm and 16mm,
    ]
    \node[box] (mains) {main-zh.tex / main-en.tex};
    \node[box, below=4mm of mains] (secs) {sections/};
    \node[box, below=4mm of secs] (figs) {figures/};
    \node[box, below=4mm of figs] (bib) {refs.bib};

    \node[core, right=18mm of secs, yshift=-4mm] (cls)
      {\textbf{boogonreport.cls}\\[2pt]{\scriptsize 文档类}};

    \node[outbox, right=16mm of cls, yshift=12mm] (pdf) {main-zh.pdf / main-en.pdf};
    \node[outbox, right=16mm of cls, yshift=-12mm] (ci) {GitHub Actions};

    \draw[arr] (mains.east) -- (cls.west |- mains.east);
    \draw[arr] (bib.east) -- (cls.west |- bib.east);
    \draw[arr] (cls.east) -- ++(6mm,0) |- (pdf.west);
    \draw[arr] (cls.east) -- ++(6mm,0) |- (ci.west);

    \node[note, below=6mm of cls] {构建：build.ps1 或 latexmk\\英文 pdflatex，中文 xelatex};
  \end{tikzpicture}
  \caption{仓库结构与构建链路。内容与类文件分离：换项目时只复制
  \texttt{boogonreport.cls} 与自己的 \texttt{sections/}、\texttt{figures/}，
  构建脚本与 CI 保持不变。}
  \label{fig:arch}
\end{figure}


数字的写法比字体更容易出问题，因为错误不会报错，只会让读者困惑。表~\ref{tab:numbers}
是我们在项目里实际踩过的几处。

\begin{table}[htbp]
  \centering
  \small
  \caption{数字写法约定。中间是排版后的样子，右边是应当写进源码的形式。}
  \label{tab:numbers}
  \begin{tabular}{@{}lll@{}}
    \toprule
    场景 & 排版结果 & 源码写法 \\
    \midrule
    千分位        & 32,768      & \texttt{32\{,\}768} \\
    乘号          & 8$\times$H100 & \texttt{8\$\textbackslash times\$H100} \\
    数值与单位    & 38\,ms      & \texttt{38\textbackslash,ms} \\
    百分号        & 62.4\%      & \texttt{62.4\textbackslash\%} \\
    精度统一      & 62.4 与 61.0 & 同一张表不要混用 62.40 \\
    提升幅度      & +2.1 points over the 7B baseline & 写清基线，不要只写 +2.1 \\
    \bottomrule
  \end{tabular}
\end{table}

第三行常被忽略：\texttt{38ms} 在断行时可能被拆成两行，读起来像 38 和 ms 两个
东西；\verb|38\,ms| 用一个细空格把数值和单位绑在一起，断行时不会被拆开。

\section{构建与发布}
\label{sec:build}

仓库里只有三类文件：类、内容、工程脚本。\cref{fig:arch} 已经画过它们的关系，
这里的重点是让「换一台机器」不产生额外心智负担。

\subsection{本地构建}

Windows 上用 \texttt{build.ps1}，macOS 与 Linux 上用 \texttt{Makefile}，两者
都调用 \texttt{latexmk} 完成两遍编译与参考文献。脚本会自动判断语言：英文走 pdfLaTeX，
中文走 XeLaTeX，因此不需要记「这份文件该用哪个引擎」。

\begin{lstlisting}[language=bash,
  caption={常用构建命令。给出 \texttt{-Job} 时只构建其中一份。},
  label={lst:build}]
./build.ps1                 # 两份示例都构建
./build.ps1 -Job main-zh    # 只构建中文示例
./build.ps1 -Clean          # 只清理中间文件
./build.ps1 -Watch          # 改一行编译一次
make en                     # macOS / Linux
\end{lstlisting}

\subsection{持续集成}

CI 做两件事：证明干净机器能编译，以及把 PDF 留档。两个 job 并行跑，英文用
pdfLaTeX、中文用 XeLaTeX 加 Fandol 字体集；任何 push 都会编译，并把成品作为
构建产物（artifact）上传。打 tag 时额外生成一个 Release，附上两份 PDF。

\begin{lstlisting}[language=yaml,
  caption={CI 的关键片段。中文任务通过 \texttt{\textbackslash PassOptionsToClass} 注入 \texttt{fandol} 字体集，避免依赖宿主系统字体。},
  label={lst:ci}]
jobs:
  build:
    strategy:
      matrix:
        include:
          - job: main-en
            engine: pdflatex
          - job: main-zh
            engine: xelatex
    steps:
      - uses: actions/checkout@v4
      - uses: xu-cheng/latex-action@v3
        with:
          root_file: ${{ matrix.job }}.tex
\end{lstlisting}

CI 能保证的只是「编译通过」，它不知道表格有没有溢出、图有没有孤零零占一页。
排版仍然需要人眼过一遍，这一条被写进了发布检查清单。

\subsection{版本}

版本号写进 \verb|\reportsetup{version={...}}|，与 git tag 保持一致。改版时
同时更新 PDF 目视检查、\texttt{CHANGELOG.md} 和 Release 说明——这三件事
任何一件单独做，都会在半年后变成谜。

\section{局限}
\label{sec:limit}

\textbf{篇幅。}单栏宽版心适合 8--20 页的报告；再长就该换双栏，这套模板不提供。

\textbf{会议模板。}NeurIPS、ICLR、ACL 的格式以官方 \texttt{.sty} 为准，本模板不内置，
也不应该内置：再分发涉及许可问题。

\textbf{数学。}只补了五个容易静默丢失的符号（≥ ≤ ≠ ≈ ∞）；\verb|\mathbb{R}| 与
\verb|\bm| 仍需自行加载宏包。

\textbf{中文边界。}断行与避头尾交给 \texttt{ctex}，繁体与日韩文未测试。

\textbf{可复现与可复制。}CI 只验证编译、不锁定 TeX Live 版本；Fandol 字体集生成的
PDF 里中文无法检索或复制，需要可复制 PDF 时改用系统字体集。


%% ----------------------------- 参考文献 ------------------------------------
\bibliography{refs}

%% ----------------------------- 附录 ----------------------------------------
\appendix
\section{选项与命令速查}
\label{app:cheatsheet}

\begin{table}[htbp]
  \centering
  \small
  \caption{类选项。写在 \texttt{\textbackslash documentclass[...]\{boogonreport\}} 的方括号里。}
  \label{tab:options}
  \begin{tabular}{@{}ll@{}}
    \toprule
    选项 & 作用 \\
    \midrule
    \texttt{zh}          & 中文报告：切字体管线、题注与交叉引用文案（需 XeLaTeX）\\
    \texttt{logo} / \texttt{nologo} & 首屏是否显示品牌带（默认显示）\\
    \texttt{times}       & 正文与数学换成 Times 家族 \\
    \texttt{numeric}     & 参考文献改为数字引用 \texttt{[1]} \\
    \texttt{nonumbering} & 章节不编号，适合白皮书 \\
    \texttt{leftalign}   & 标题块左对齐，默认居中 \\
    \texttt{draft}       & 行号 + \verb|\todo{}| 可见 \\
    \texttt{fandol}      & 中文用 Fandol 字体集（跨平台可编译）\\
    \texttt{winfonts}    & 中文用 Windows 系统字体 \\
    \bottomrule
  \end{tabular}
\end{table}

\begin{table}[htbp]
  \centering
  \small
  \caption{常用命令与环境。}
  \label{tab:commands}
  \begin{tabular}{@{}ll@{}}
    \toprule
    命令 & 用途 \\
    \midrule
    \verb|\reportsetup{...}| & 标题、作者、编号、版本、链接与文字标 \\
    \verb|\keywords{...}|    & 关键词，放在摘要内或摘要后都能正常排版 \\
    \verb|\secref{...}|      & 中文输出「第 3 节」，英文输出「section 3」\\
    \verb|\cref{...}|        & 图表引用，中文输出「图 1」「表 2」\\
    \verb|\appref{...}|      & 附录引用，输出「附录 A」\\
    \verb|\best{}| / \verb|\second{}| & 最优值加粗 / 次优值下划线 \\
    \texttt{contributions}   & 紧凑的贡献点列表 \\
    \verb|\todo{}| / \verb|\fixme{}| & 草稿标记，仅 \texttt{draft} 下可见 \\
    \bottomrule
  \end{tabular}
\end{table}

\section{作者与致谢}
\label{app:authors}

\textbf{维护。}BoogonClothman 设计并维护本模板，包括类文件、示例、文档与构建
流程。

\textbf{致谢。}模板的版式语言受益于公开技术报告的长期实践，尤其是 DeepSeek 与
Qwen 系列报告对「一页讲清一件事」的排版处理；本书面语言与字体选择则参照
Bringhurst 与 \LaTeX{} 社区的通行做法\citep{bringhurst2013elements,mittelbach2023companion}。

\textbf{许可。}本仓库以 MIT 许可发布，可商用、可修改、可再分发。字体随
TeX Live 分发，按各自的 LPPL/GPL 类许可嵌入 PDF。

\textbf{获取。}\url{https://github.com/BoogonClothman/boogon-techreport-kit}


\end{document}
